\chapter{Tattoo}
\index{Tattoo}

\medskip
\noindent\textit{Educational reference; always seek professional advice for individual care.}

\section*{Definition and Scope}
A tattoo is a permanent form of body art made by injecting coloured ink into the deeper layers of the skin using fine needles. The word "tattoo" is also used for the finished design, the process of creating it, and the practice as a whole. Tattoos can be decorative, cultural, or commemorative, but they also have medical and cosmetic uses, such as reconstructing the look of a nipple after breast surgery, covering scars, or tattooing eyebrows, lips, and eyeliner. There is also temporary body art, such as henna and stick-on transfers, which sit on the surface of the skin and fade over time. In many countries, tattooing is a regulated industry: in most states and territories, tattooists must be registered and premises must meet health standards, because the process involves breaking the skin with equipment that must be sterile. Tattoos are very common and, for most healthy people, safe when performed by a licensed professional in a clean studio and cared for properly afterwards. They have become steadily more popular and accepted across all age groups in recent decades, and most people either have a tattoo or know someone who does. This chapter explains how tattoos are done, what to expect during and after the process, the benefits and the risks, and the practical steps that reduce the chance of infection, scarring, or regret. It covers both decorative tattoos and the medical and cosmetic uses of the technique, and it is a general guide rather than a substitute for advice from a doctor, a dermatologist, or a licensed tattooist.

\section*{Benefits and Purpose}
People get tattoos for many personal reasons, and the benefits are as much emotional and social as they are cosmetic. For many, a tattoo is a form of self-expression that carries a story: a memorial to a loved one, a symbol of a milestone or a belief, an expression of identity, or a mark of belonging to a family, group, or culture. In many Aboriginal and Torres Strait Islander communities and in Pacific cultures, tattooing has deep ceremonial and cultural meaning that stretches back generations. Other people use tattoos to reclaim their bodies, for example by covering scars from surgery, injury, or self-harm, or to restore appearance after a mastectomy, where tattooed areola reconstruction can be a genuinely therapeutic step in recovery. Some people simply enjoy the aesthetic of wearing art, and the process itself, which often involves a meaningful design and time with an artist, can be a positive experience. There is also evidence that tattoos give many people a sense of confidence and control over their own bodies. The benefits depend heavily on the decision being voluntary, well-considered, and carried out safely, which is why preparation and choosing the right artist matter as much as the design itself. For the same reasons, professionals who tattoo medical areas such as areola reconstruction receive specific training, and their work is done in close consultation with the treating health team.

\section*{How It Works}
Tattooing works by placing ink permanently in the dermis, the layer of skin below the surface layer (the epidermis). Because the epidermis constantly sheds and renews, ink placed only in the top layer would disappear quickly; the artist's needles therefore push ink to a depth of roughly one to two millimetres, below the surface but not into the fat layer beneath the skin. The tattoo machine holds a set of fine needles that puncture the skin very rapidly, usually between fifty and three thousand times per minute, and each puncture deposits a tiny drop of ink. Modern machines are driven by electricity, either with a rotating coil mechanism or a linear motor, and the artist can vary the needle depth, speed, and grouping to create fine lines, shading, or solid colour. Before work begins, the design is stencilled onto the skin, the area is shaved and cleaned with antiseptic, and the artist wears gloves and uses a fresh, sterile, single-use needle and ink cartridge for every client. A numbing cream or spray is sometimes applied to reduce discomfort. The finished tattoo is cleaned, an antibacterial ointment may be applied, and the area is bandaged to protect it as it begins to heal. The healing process then proceeds in stages, from an open wound, to a scabbed surface, to the new skin that finally holds the ink in place.

\section*{What to Expect}
Getting a tattoo takes time, and the experience is different for everyone. The process itself causes some pain, which most people describe as a stinging or scratching sensation; how much it hurts depends on the site, the size of the design, and the person's own tolerance. Areas over bone, such as the ribs, spine, and shins, and areas with thin skin, such as the inside of the wrist and the ankle, tend to be more sensitive, while fleshy areas such as the upper arm and thigh are usually more comfortable. A small tattoo can take less than an hour, while a large, detailed piece can take several hours or several sessions spread over weeks. Bleeding is normal in small amounts, and the area may be sore, swollen, and red for the first day or two. For the next one to three weeks the tattoo may weep, then scab and peel as the top layers of skin repair, and during this time it will itch. It is essential not to pick at the scabs, because this can pull out ink and leave patchy or scarred results. The surface skin heals within about two to four weeks, but the deeper layers continue to settle for months, and the true colour is usually not seen fully until after the first month. Most people need a "touch-up" session later if parts of the design have faded unevenly. Tattoos are permanent; while laser removal can lighten them, it is expensive, takes many sessions, and may not remove the ink completely. It is also worth knowing that a tattoo will never look exactly like the stencil or the artist's other work: skin is living tissue, and the finished result depends on healing, skin type, and the placement chosen. A good artist will talk you through what is realistic for your skin and your design before any ink is used.

\section*{When to Seek Professional Advice}
Most tattoo aftercare is managed at home, but some situations need medical review. See a doctor if you develop signs of infection, such as spreading redness, increasing pain, warmth, swelling, discharge of pus, or a fever; if the tattoo is not healing within a few weeks; or if you develop an itchy, raised, or weeping rash that suggests an allergic reaction to the ink. People with diabetes, a weakened immune system, a bleeding disorder, or a history of keloid scarring should discuss tattoos with their doctor before getting one, as should people who are pregnant or breastfeeding. Anyone who has had a serious condition such as an artificial heart valve or who is on medicines that affect healing should also seek advice. If you are thinking of tattooing over a mole or an area of changed skin, see a doctor or dermatologist first, because the design can hide changes that need monitoring. Finally, if you are considering having a tattoo removed, discuss the options, including laser removal and surgical excision, with a practitioner experienced in tattoo removal.

\textbf{Contact your local emergency services for:}
\begin{itemize}
    \item Sudden severe pain, swelling, or a rapidly spreading area of redness and heat after a tattoo, which can suggest a serious skin infection
    \item A high fever with chills
    \item Difficulty breathing or swelling of the face and lips after a tattoo or a henna application
\end{itemize}

\section*{Risks and Cautions}
When a tattoo is done by a licensed professional with sterile equipment, the risks are low, but they are not zero, and they deserve an honest look before the decision is made.

\begin{itemize}
    \item \textbf{Infection:} any needle in the skin can introduce bacteria, causing local skin infection or, rarely, deeper infection such as a blood infection or heart infection; using a registered studio with single-use needles is the main protection
    \item \textbf{Blood-borne viruses:} hepatitis B, hepatitis C, and HIV can be transmitted if equipment is reused or not properly sterilised, which is why never accepting a "homemade" or unlicensed tattoo is critical
    \item \textbf{Allergic reactions:} some inks, particularly red, yellow, and certain white pigments, can cause allergic skin reactions that appear days or even years later, sometimes as an itchy, raised rash
    \item \textbf{Sun sensitivity:} some pigments become irritated or raised when exposed to strong sunlight, and sun exposure fades the ink unevenly, so sun protection over a tattoo is important for life
    \item \textbf{Scarring and keloids:} some people form raised scars or keloids, and tattooing over existing scars can be unpredictable
    \item \textbf{MRI interference:} some tattoo pigments can cause a warming sensation or mild artefact during magnetic resonance imaging, though serious problems are rare
    \item \textbf{Fading and distortion:} sunlight, age, and weight change fade and blur tattoos over time, and colours behave differently on different skin tones
    \item \textbf{Regret:} a significant minority of people later wish they had not been tattooed, and removal is costly, slow, and often incomplete
    \item \textbf{PPD in henna:} "black henna" temporary tattoos often contain a strong dye called para-phenylenediamine, which can cause severe allergic reactions and permanent skin damage
    \item \textbf{Moles and skin checks:} a tattoo over a mole can make it hard for a doctor to examine it, and an ink reaction can mimic or mask skin change
\end{itemize}

\section*{Tips and Practical Advice}
\begin{itemize}
    \item Choose a licensed, registered tattoo studio in your local, and check online reviews and the artist's portfolio before booking
    \item Ask about sterilisation: needles, ink, and tubes should be single-use, and the artist should open them fresh in front of you
    \item Discuss allergies, existing medical conditions, and any medicines you take with the artist and your doctor before you go
    \item Do not drink alcohol or use recreational drugs before your session, as they thin the blood, increase bleeding, and cloud judgement
    \item Have a meal before your appointment and drink plenty of water, as tattooing is tiring and can cause light-headedness
    \item Follow the aftercare instructions exactly: keep the bandage on as advised, wash gently with plain soap and water, and apply only the recommended moisturiser
    \item Never scratch or pick a healing tattoo, and keep it out of direct sunlight, pools, and salt water until it is fully healed
    \item Never get a tattoo from an unlicensed or "kitchen table" operation, and never share equipment
    \item If you are a parent, respect your state's rules on age limits, and support younger people in waiting until they are old enough and sure
    \item Take your time choosing a design and artist, and sleep on the decision before booking, since tattoos are permanent
\end{itemize}

\section*{Sources and Further Reading}
The following organisations provide reliable information about tattoos, tattoo safety, and skin health.

\begin{itemize}
    \item healthdirect Australia. \textit{Tattoos: safety, aftercare, and risks.} https://www.healthdirect.gov.au/
    \item Australian Department of Health and Aged Care. \textit{Tattooing and body art safety guidance.} https://www.health.gov.au/
    \item NHS. \textit{Tattoos: risks, aftercare, and when to seek help.} https://www.nhs.uk/
    \item Mayo Clinic. \textit{Tattoos: what to know before you get one.} https://www.mayoclinic.org/
    \item Cancer Council Australia. \textit{Skin health and tattoo considerations.} https://www.cancer.org.au/
    \item World Health Organization (WHO). \textit{Blood safety and body art practices.} https://www.who.int/
\end{itemize}

\clearpage
